\documentclass[11pt,a4paper]{article}
\usepackage{gtu-el-yazisi}
\begin{document}
\notbilgi{Sayısal İşaret İşleme}{Fark denklemleri, FIR/IIR, DTFT ve frekans cevabı}{Dsp-ders-3.pdf}

\sayfa{1}
\elcizimi[0.4]{sekiller/Dsp-ders-3-s1-1.png}[1]
FIR : sonlu sayıda 0'dan farklı değer alır

\elcizimi[0.5]{sekiller/Dsp-ders-3-s1-2.png}[1]
ve $x[n]=\sum_{k=-\infty}^{\infty} x[k]\,\delta[n-k]$

\begin{align*}
y[n] &= x[n]*h[n]\\
y[n] &= \sum_{k=-\infty}^{\infty} x[k]\,h[n-k]\\
\intertext{\small $\downarrow$ \okunamadi{bunse} \qquad $\swarrow$ \okunamadi{bu dummy}}
&= \sum_{k=-\infty}^{\infty} h[k]\,x[n-k]
\end{align*}

\elcizimi[0.5]{sekiller/Dsp-ders-3-s1-3.png}[1]
\elcizimi[0.55]{sekiller/Dsp-ders-3-s1-4.png}[1]
$\rightarrow$ olsa bir yerden sonra toplamı kesebiliriz

Filtre IIR ise konvülasyon sadece kağıt üzerinde gerçekleşebilir\\
yani toplam formülü kullanılamaz.

\begin{align*}
y[n] + \sum_{k=1}^{N} a_k\,y[n-k] &= \sum_{k=0}^{M} b_k\,x[n-k]\\
y[n] &= \sum_{k=0}^{M} b_k\,x[n-k] - \sum_{k=1}^{N} a_k\,y[n-k]
\end{align*}
\vurgu[turuncu]{Constant coefficient Difference Equation}

\begin{align*}
y[n] &= b_0x[n]+b_1x[n-1]+b_2x[n-2]+\cdots+b_Mx[n-M]\\
&\quad -a_1y[n-1]-a_2y[n-2]
\end{align*}

Sabit katsayılı fark denklemi

Bu nereden geliyor??

Bir sistem düşünelim\footnote{Orijinalde turuncu kalemle üstü çizilen terimler turuncu vurguyla gösterilmiştir.}
\[ y[n] = \tfrac{1}{3}x[n] + \text{\vurgu[turuncu]{$\tfrac{1}{3}x[n-1]$}} + \text{\vurgu[turuncu]{$\tfrac{1}{3}x[n-2]$}} \]
\[ \longrightarrow\ h[n] = \tfrac{1}{3}\delta[n] + \tfrac{1}{3}\delta[n-1] + \tfrac{1}{3}\delta[n-2] \]
\[ y[n-1] = \text{\vurgu[turuncu]{$\tfrac{1}{3}x[n-1]$}} + \text{\vurgu[turuncu]{$\tfrac{1}{3}x[n-2]$}} + \tfrac{1}{3}x[n-3] \]

\begin{align*}
y[n]-y[n-1] &= \tfrac{1}{3}x[n]-\tfrac{1}{3}x[n-3]\\
y[n] &= \tfrac{1}{3}x[n]-\tfrac{1}{3}x[n-3]+y[n-1]
\end{align*}

\begin{align*}
b_0 &= \tfrac{1}{3} & a_1 &= 1\footnotemark\\
b_1 &= 0 & a_2 &= 0\\
b_2 &= 0 & a_3 &= 0\\
b_3 &= -\tfrac{1}{3} & &\\
b_4 &= 0 & &
\end{align*}
\footnotetext{Claude notu: yukarıdaki fark denklemine göre muhtemelen $a_1=-1$ olmalı.}

\sayfa{2}
\section{Örnek}
\[ y[n]=\tfrac{1}{100}x[n]+\tfrac{1}{100}x[n-1]+\cdots+\tfrac{1}{100}x[n-99] \]
\begin{align*}
y[n]-y[n-1] &= \tfrac{1}{100}x[n]-\tfrac{1}{100}x[n-100]\\
y[n] &= \tfrac{1}{100}x[n]-\tfrac{1}{100}x[n-100]+y[n-1]
\end{align*}

\[ y[n]=5x[n]+3 \qquad \Rightarrow\ \begin{array}{l}\text{lineer değil}\\ \text{LTI değil}\end{array} \]

\elcizimi[0.5]{sekiller/Dsp-ders-3-s2-1.png}[2]
\textbf{Başlangıç koşulları}
\begin{align*}
y[-4] &= 0 & y[-6] &= 0\\
y[-5] &= 0 & y[-1] &= y[0] = 0
\end{align*}

\section{DTFT}
$\rightarrow$ sonlu olmalı
\[ X(e^{j\omega})=\sum_{n=-\infty}^{\infty} x[n]\,e^{-j\omega n} \]
$\omega$'nın fonksiyonu sürekli
\elcizimi[0.6]{sekiller/Dsp-ders-3-s2-2.png}[2]

\begin{gather*}
\cos(\Omega t)\Big|_{t=nT}\\
t=\frac{n}{F_s}\\
\cos\Big(\frac{\Omega}{F_s}\,n\Big)\\
\omega=\frac{\Omega}{F_s}\\
\cos(\omega n)\\
-\infty<\Omega<\infty\\
-\infty<\omega<\infty
\end{gather*}

$X(e^{j\omega})$ $\omega$'nın $2\pi$'sinde periyodik
\begin{gather*}
0\le\omega<2\pi\\
-\pi\le\omega<\pi
\end{gather*}
(ses işareti 3,2 | 4 kHz'e sınırlanıyor\\
8 kHz örnekleme frekansı)

\begin{align*}
X(e^{j\omega}) &= X_R(e^{j\omega})+jX_I(e^{j\omega})\\
&= |X(e^{j\omega})|\,e^{j\angle X(e^{j\omega})}
\end{align*}
$\nearrow$ Fourier için yeter koşuldur
\begin{align*}
|X(e^{j\omega})| &= \Big|\sum_{k=-\infty}^{\infty} x[k]\,e^{-j\omega k}\Big|
\le \sum_{k=-\infty}^{\infty} \big|x[k]\,e^{-j\omega k}\big|\\
&\le \sum_{k=-\infty}^{\infty} |x[k]| \le \infty
\end{align*}

\sayfa{3}
$\downarrow$ bu şartı sağlamayan\\
FIR sistemlerin her zaman fourieri\\
IIR sistemlerde sadece kararlıysa var

\section{Properties of DTFT}
\[ X(e^{j\omega})=X(e^{j\omega}e^{j2\pi})=X(e^{j(\omega+2\pi)}) \qquad \text{Periyodik } 2\pi \text{ ile} \]

\begin{align*}
x_1[n] &\longleftrightarrow X_1(e^{j\omega})\\
x_2[n] &\longleftrightarrow X_2(e^{j\omega})\\
ax_1[n]+bx_2[n] &\longleftrightarrow aX_1(e^{j\omega})+bX_2(e^{j\omega})
\end{align*}
\[ \cos(\omega n)=\tfrac{1}{2}e^{j\omega n}+\tfrac{1}{2}e^{-j\omega n} \]

\begin{align*}
x[n] &\overset{\mathcal{F}}{\longleftrightarrow} X(e^{j\omega})\\
x[n-n_0] &\overset{\mathcal{F}}{\longleftrightarrow} e^{-j\omega n_0}X(e^{j\omega})\\
e^{j\omega_0 n}x[n] &\overset{\mathcal{F}}{\longleftrightarrow} X(e^{j(\omega-\omega_0)})
\end{align*}

\vurgu[turuncu]{ispat}
\[
\left\{\begin{aligned}
X(e^{j\omega}) &= \sum_{n=-\infty}^{\infty} x[n]\,e^{-j\omega n}\\
X_1(e^{j\omega}) &= \sum_{n=-\infty}^{\infty} e^{j\omega_0 n}\,x[n]\,e^{-j\omega n}\\
&= \sum_{n=-\infty}^{\infty} x[n]\,e^{j(\omega-\omega_0)n}\footnotemark
\end{aligned}\right.
\]
\footnotetext{Claude notu: muhtemelen $e^{-j(\omega-\omega_0)n}$ olmalı.}

\[ \omega_0=\pi \qquad e^{j\pi n}=(-1)^n \]
\begin{align*}
h[n] &\longleftrightarrow H(e^{j\omega})\\
(-1)^n h[n] &\longleftrightarrow H(e^{j(\omega-\pi)})
\end{align*}

\elcizimi[0.7]{sekiller/Dsp-ders-3-s3-1.png}[3]
alçak geçiren

yüksek geçiren çünkü en yüksek frekans $\pi$

\begin{align*}
x[n] &\longleftrightarrow X(e^{j\omega})\\
x^{*}[n] &\longleftrightarrow X^{*}(e^{-j\omega}) \qquad \leftarrow \text{konjüge simetrik}
\end{align*}
\[ x[n]\ \text{reel} \ \Rightarrow\ x[n]=x^{*}[n] \]
\[ X(e^{j\omega})=X^{*}(e^{-j\omega}) \qquad \Rightarrow\ |X(e^{j\omega})|=|X^{*}(e^{-j\omega})| \]
$\nearrow$ Fazlar orjine göre simetriktir

$\ast$ Reel bir işaretin fourier dönüşümü $\omega$'ya simetriktir

\sayfa{4}
\begin{align*}
x[n] &\longleftrightarrow X(e^{j\omega})\\
x[-n] &\longleftrightarrow X(e^{-j\omega})
\end{align*}

\vurgu[turuncu]{ispat}
\[
\left\{\begin{aligned}
X(e^{j\omega}) &= \sum_{n=-\infty}^{\infty} x[n]\,e^{-j\omega n}\\
X_1(e^{j\omega}) &= \sum_{n=-\infty}^{\infty} x[-n]\,e^{-j\omega n}\\
X_1(e^{j\omega}) &= \sum_{n=-\infty} x[k]\,e^{j\omega k}\\
&\hspace{6em}\downarrow\ e^{-j(-\omega)k}
\end{aligned}\right.
\]

\begin{align*}
x[n]-x[n-1] &\overset{\mathcal{F}}{\longleftrightarrow} X(e^{j\omega})-e^{-j\omega}X(e^{j\omega})\\
&= (1-e^{-j\omega})\,X(e^{j\omega})
\end{align*}
\okunamadi{türev alabilir}

\elcizimi[0.55]{sekiller/Dsp-ders-3-s4-1.png}[4]
yüksek frekanslar güçlenir

\[ n\,x[n] \longleftrightarrow j\,\frac{dX(e^{j\omega})}{d\omega} \]
\[ \sum_{n=-\infty}^{\infty}|x[n]|^{2} = \frac{1}{2\pi}\int_{2\pi}|X(e^{j\omega})|^{2}\,d\omega \]
\vurgu[turuncu]{x[n] işaretinin enerjisi} \qquad \vurgu[turuncu]{Enerji yoğunluğu}

\elcizimi[0.45]{sekiller/Dsp-ders-3-s4-2.png}[4]

\section{The convolution Theorem}
\begin{align*}
\text{\textcircled{1}}\quad x_1[n]*x_2[n] &\overset{\mathcal{F}}{\longleftrightarrow} X_1(e^{j\omega})\,X_2(e^{j\omega})\\
\text{\textcircled{2}}\quad x_1[n]\,x_2[n] &\overset{\mathcal{F}}{\longleftrightarrow} \tfrac{1}{2\pi}\,X_1(e^{j\omega})\,X_2(e^{j\omega})\footnotemark\\
y[n]=x[n]*h[n] &\overset{\mathcal{F}}{\longleftrightarrow} Y(e^{j\omega})=X(e^{j\omega})H(e^{j\omega})
\end{align*}
\footnotetext{Claude notu: muhtemelen $X_1(e^{j\omega})$ ile $X_2(e^{j\omega})$ arasında periyodik konvolüsyon ($\circledast$) olmalı.}

\elcizimi[0.55]{sekiller/Dsp-ders-3-s4-3.png}[4]
\elcizimi[0.38]{sekiller/Dsp-ders-3-s4-4.png}[4]

\sayfa{5}
\elcizimi[0.58]{sekiller/Dsp-ders-3-s5-1.png}[5]
\begin{align*}
y[n] &= \sum_{k=-\infty}^{\infty} h[k]\,x[n-k]\\
&= \sum_{k=-\infty}^{\infty} h[k]\,e^{j\omega(n-k)}\\
&= \Big[\sum_{k=-\infty}^{\infty} h[k]\,e^{-j\omega k}\Big]e^{j\omega n}\\
&= e^{j\omega n}\,H(e^{j\omega})
\end{align*}

özvektör
\elcizimi[0.3]{sekiller/Dsp-ders-3-s5-2.png}[5]
\[ A\vec{x}=\lambda\vec{x} \]
eigen value (öz değer) \qquad eigen vector (öz vektör)

\elcizimi[0.4]{sekiller/Dsp-ders-3-s5-3.png}[5]
\[ y[n]=e^{j\frac{\pi}{2}n}\,H\big(e^{j\frac{\pi}{2}}\big) \]
\vurgu[turuncu]{dizi} \quad \vurgu[turuncu]{katsayı} \quad \vurgu[turuncu]{(örnek $0{,}4\,e^{j\theta}$)}

$\omega=\frac{\pi}{2}$ giriş için

\vurgu[turuncu]{frekans değişmiyor}

Lineer bir sistem giriş işaretin frekansı korunur.\\
Fakat 0'a çevirebilir.

\elcizimi[0.4]{sekiller/Dsp-ders-3-s5-4.png}[5]

Girişin frekansına göre sistem farklı davranır bu yüzden $H(e^{j\omega})$'ya frekans cevabı denir.

\elcizimi[0.52]{sekiller/Dsp-ders-3-s5-5.png}[5]

\sayfa{6}
\begin{align*}
H(e^{j\omega}) &= \sum_{n=-\infty}^{\infty} h[n]\,e^{-j\omega n}\\
&= \sum_{n=-\infty}^{\infty} \big(\tfrac{1}{2}\big)^{n}u[n]\,e^{-j\omega n}\\
&= \sum_{n=0}^{\infty} \Big(\underbrace{\tfrac{1}{2}e^{-j\omega}}_{\alpha}\Big)^{n} = \frac{1}{1-\frac{1}{2}e^{-j\omega}}
\end{align*}

\elcizimi[0.42]{sekiller/Dsp-ders-3-s6-1.png}[6]
\elcizimi[0.45]{sekiller/Dsp-ders-3-s6-2.png}[6]

\begin{align*}
y[n] &= e^{j\frac{\pi}{2}n}\,\underbrace{H(e^{j\frac{\pi}{2}})}_{0{,}89\,e^{-j\,0{,}4636}}\\
&= 0{,}89\,e^{j\left(\frac{\pi}{2}n-0{,}4636\right)}
\end{align*}

\end{document}
